\documentclass[11pt]{article}
\usepackage[margin=1in]{geometry}
\usepackage{amsmath,amssymb,amsthm,hyperref}
\newtheorem{theorem}{Theorem}
\newtheorem{lemma}[theorem]{Lemma}
\title{A sharp order bound for an integral of the modified Bessel function}
\author{Principia Math In Progress}
\date{11 September 2026}
\begin{document}
\maketitle
\begin{abstract}
We prove Gaunt's conjectured lower bound for an integral of the modified
Bessel function of the second kind, including the boundary at which the
optimal order is zero. The bound is sharp among all fixed real orders.
The key observation is strict convexity of $s\mapsto K_{\sqrt{s}}(x)$.
\end{abstract}

\section*{Task 1: Resolve MathDB problem 339406}
Conjecture 2.4, equation (2.17), page 7 of \cite{gaunt}, is the following
theorem. We also prove its assertion about the best possible order.

\begin{theorem}\label{thm:main}
Let $n<1$ and $\nu\geq1-n+\sqrt{2(1-n)}$. Put
\[
 m=\nu+n,\qquad d=\sqrt{m^2-2\nu-1},\qquad
 I(x)=\int_x^\infty t^{-\nu}K_m(t)\,dt.
\]
Then, for every $x>0$,
\begin{equation}\label{eq:bound}
 I(x)>x^{-\nu}K_d(x).
\end{equation}
For any real $\beta$, put $\rho=|m-\beta|$. If $\rho<d$, then
$K_\rho(x)<K_d(x)$ for every $x>0$. If $\rho>d$, then
$I(x)<x^{-\nu}K_\rho(x)$ for all sufficiently large $x$.
Thus the optimal shifts are exactly $\beta=m-d$ and $\beta=m+d$,
which give the same function and coincide when $d=0$.
For $n=1$ the corresponding identity, valid for every real $\nu$, is
\begin{equation}\label{eq:equality}
 \int_x^\infty t^{-\nu}K_{\nu+1}(t)\,dt=x^{-\nu}K_\nu(x).
\end{equation}
\end{theorem}

\subsection*{Convexity in the squared order}
We use the standard real-order representation
\begin{equation}\label{eq:representation}
 K_a(x)=\int_0^\infty e^{-x\cosh u}\cosh(au)\,du,
 \qquad a\in\mathbb R,\quad x>0.
\end{equation}
The recurrence and derivative identities used below are also recorded
in Appendix B of \cite{gaunt}.

\begin{lemma}\label{lem:convex}
For fixed $x>0$, the function $F_x(s)=K_{\sqrt{s}}(x)$ is strictly
increasing and strictly convex on $[0,\infty)$. Moreover,
\begin{equation}\label{eq:tangent}
 F_x'(1)=\frac{K_0(x)}{2x}.
\end{equation}
\end{lemma}
\begin{proof}
For $u>0$,
\[
 \cosh(u\sqrt{s})=\sum_{j=0}^\infty\frac{u^{2j}s^j}{(2j)!}
\]
is strictly increasing and strictly convex on $s\geq0$; its first and
second derivatives are positive, including at zero. Integration against
the positive weight in \eqref{eq:representation} preserves both strict
inequalities. Differentiation is justified locally uniformly in $s$:
on compact order intervals, every required order derivative is dominated
by a constant times a polynomial in $u$ times
$e^{-x\cosh u+A u}$ for some fixed $A$, an integrable function.
Finally, integration by parts gives
\[
 2F_x'(1)=\int_0^\infty u\sinh u\,e^{-x\cosh u}\,du
 =\frac1x\int_0^\infty e^{-x\cosh u}\,du=\frac{K_0(x)}x.
\]
The boundary term $u e^{-x\cosh u}$ vanishes at both endpoints.
\end{proof}

\subsection*{The derivative comparison}
The assumptions imply
\begin{equation}\label{eq:domain}
 d^2=(m-1)^2-2(1-n),\qquad 0\leq d<m-1<\nu.
\end{equation}
Indeed, $m-1\geq\sqrt{2(1-n)}>0$ and $n<1$.
In particular,
\[
 m^2-(d+1)^2=2(\nu-d)>0.
\]
First suppose $d>0$. Apply strict convexity of $F_x$ to the three ordered
points $a=(d-1)^2<b=(d+1)^2<c=m^2$. The strictly increasing secant slopes give
\begin{align*}
 K_m(x)
 &>K_{d+1}(x)+\frac{\nu-d}{2d}
             \bigl(K_{d+1}(x)-K_{d-1}(x)\bigr)\\
 &=\frac{\nu+d}{2d}K_{d+1}(x)
    -\frac{\nu-d}{2d}K_{d-1}(x)\\
 &=K_{d-1}(x)+\frac{\nu+d}{x}K_d(x),
\end{align*}
where $K_{-a}=K_a$ and
$K_{d+1}-K_{d-1}=2dK_d/x$ were used.
If $d=0$, then $m^2=1+2\nu>1$. Strict convexity and
\eqref{eq:tangent} instead give
\[
 K_m(x)>F_x(1)+2\nu F_x'(1)
       =K_1(x)+\frac\nu x K_0(x).
\]
Consequently, in the full parameter range, including $d=0$,
\begin{equation}\label{eq:comparison}
 K_m(x)>K_{d-1}(x)+\frac{\nu+d}{x}K_d(x).
\end{equation}
This proves the derivative inequality proposed in equation (2.18) of
\cite{gaunt}; at $d=0$ it gives its nonsingular form.

\subsection*{Proof of the integral bound and sharpness}
Let $D(x)=I(x)-x^{-\nu}K_d(x)$. The identity
$K_d'=-K_{d-1}-dK_d/x$ and \eqref{eq:comparison} yield
\[
 D'(x)=x^{-\nu}\left(K_{d-1}(x)+\frac{\nu+d}{x}K_d(x)-K_m(x)\right)<0.
\]
Both terms defining $D$ tend to zero at infinity, by exponential decay
of the Bessel functions. Thus
\[
 D(x)=\int_x^\infty t^{-\nu}
 \left(K_m(t)-K_{d-1}(t)-\frac{\nu+d}{t}K_d(t)\right)dt>0,
\]
proving \eqref{eq:bound} with strictness.

For completeness we supply the optimality calculation suggested in the
source. For each fixed real $a$, the standard expansion is
\[
 K_a(x)=\sqrt{\frac\pi2}e^{-x}x^{-1/2}
 \left(1+\frac{4a^2-1}{8x}+O(x^{-2})\right).
\]
Integration by parts for the first two terms, with the same order of
remainder after integration, gives
\begin{align*}
 I(x)&=\sqrt{\frac\pi2}e^{-x}x^{-\nu-1/2}
 \left(1+\frac{(4m^2-1)/8-\nu-1/2}{x}+O(x^{-2})\right)\\
 &=\sqrt{\frac\pi2}e^{-x}x^{-\nu-1/2}
 \left(1+\frac{4d^2-1}{8x}+O(x^{-2})\right).
\end{align*}
It follows that
\begin{equation}\label{eq:sharp}
 I(x)-x^{-\nu}K_\rho(x)
 =\sqrt{\frac\pi2}e^{-x}x^{-\nu-1/2}
   \left(\frac{d^2-\rho^2}{2x}+O(x^{-2})\right).
\end{equation}
If $\rho>d$, the expression is eventually negative. If $0\leq\rho<d$,
Lemma \ref{lem:convex} gives $K_\rho<K_d$ pointwise. The only remaining
orders are $\rho=d$, equivalent to the two stated shifts.

Finally, for $n=1$, the identity
$(x^{-\nu}K_\nu(x))'=-x^{-\nu}K_{\nu+1}(x)$ proves
\eqref{eq:equality} after integration to infinity. It holds for all
real $\nu$, without imposing the $n<1$ parameter condition.

\section*{Task 2: Classify the parameter range for sharp order bounds}
Determine all real pairs $(\nu,n)$ for which a fixed real order $r$
gives $\int_x^\infty t^{-\nu}K_{\nu+n}(t)\,dt\geq x^{-\nu}K_r(x)$
for every $x>0$. Identify when the large-$x$ obstruction
$r^2\leq(\nu+n)^2-2\nu-1$ is sufficient outside the range proved here.

\section*{Task 3: Quantify the positive remainder}
Use the strictly positive convexity remainder in
\eqref{eq:comparison} to obtain explicit two-sided bounds for $D(x)$.
Seek estimates that remain useful as $d\downarrow0$ and $n\uparrow1$,
and that recover the first nonzero large-$x$ term.

\section*{Task 4: Extend the sharp bound to exponential weights}
Determine optimal Bessel orders and prefactors in bounds for
$\int_x^\infty e^{\gamma t}t^{-\nu}K_{\nu+n}(t)\,dt$ when $\gamma<1$.
Establish which weighted analogues admit a convexity proof.

\section*{Task 5: Develop a general squared-order comparison principle}
Identify other special-function families admitting a positive
$\cosh(\sqrt{s}\,u)$ representation and useful order recurrences.
Derive sharp integral comparisons from their secant and tangent
inequalities, including parameter values where the secant degenerates.

\begin{thebibliography}{9}
\bibitem{gaunt} R. E. Gaunt,
\emph{Inequalities for some integrals involving modified Bessel functions},
Proc. Amer. Math. Soc. \textbf{147} (2019), 2937--2951.
\href{https://doi.org/10.1090/proc/14433}{doi:10.1090/proc/14433}.
Author-linked manuscript:
\href{https://arxiv.org/abs/1806.00524}{arXiv:1806.00524v2} (9 October 2018),
Conjecture 2.4 and Remark 2.5, page 7; standard identities in Appendix B.
\end{thebibliography}
\end{document}
